% SOURCE_AUTHORITY: sga5_fr_workpass.tex, lines 11268--11289 (LF-normalized unit).
% SOURCE_SHA256: 791F4EFFC5E02832D5D77ED03518C8156D6F07E4C8238B03545DB93D883FBB28
Pasemos ahora a la «fórmula clave» mencionada en la introducción. Las notaciones adaptadas, recordadas en su mayor parte, son las de 4. Se observará que son compatibles con las introducidas en el presente apartado, salvo que el esquema denotado por $Y'$ en 4 se denota aquí por $H$ (por «hiperplano»). Utilizaremos, por lo demás, ambas notaciones indistintamente.

\begin{statement}{Proposición 9.6}
Sean $k$ un cuerpo separablemente cerrado e $i:Y\hookrightarrow X$ una inmersión cerrada, de codimensión pura $d$, de $k$-esquemas cuasiproyectivos y lisos. Denotando por $f:X'\to X$ el $X$-esquema obtenido por la explosión de $X$ a lo largo de $Y$, sea
\[
\begin{tikzcd}
Y'\arrow[r,"j"]\arrow[d,"g"']&X'\arrow[d,"f"]\\
Y\arrow[r,"i"']&X
\end{tikzcd}
\]
el diagrama cartesiano construido a partir de $f$ e $i$. Entonces, denotando por $F$ el $\mathcal O_{Y'}$-módulo localmente libre definido por la sucesión exacta canónica
\[
0\to F\to g^*(N)\to\mathcal O_{Y'}(1)\to0,
\tag{9.6.1}
\]
se tiene la igualdad
\[
f^*i_*(y)=j_*(g^*(y)c_{d-1}(\check F))
\tag{9.6.2}
\]
en $C^*(X')$, para todo $y\in C^*(Y)$.
\end{statement}
